\documentclass{article}
\usepackage[margin=1in]{geometry}
\usepackage{problemset}
\showsolntrue

% Custom commands
\newcommand{\R}{\mathbb{R}}
\newcommand{\X}{\mathcal{X}}
\newcommand{\Y}{\mathcal{Y}}

\begin{document}

\pstitle{%
    CSE 151A%
}{
    Homework 8%
}[due Thursday, May 29 at 11:59PM]

\textbf{Instructions:}
\begin{itemize}
    \item You may discuss problems with your study group, but ultimately all your work (mathematical problems, code, experimental details) must be individual.
    \item Your solutions must be \texttt{typed up} and uploaded to Gradescope by 11.59PM on Thursday May 29. No late homeworks will be accepted under any circumstances, so you are encouraged to upload early.
    \item A subset of the problems will be graded.
\end{itemize}

\begin{probset}
\begin{prob}
    A linear predictor is used to solve a classification problem with three classes. The data is two-dimensional and the linear functions for each class are:

    \begin{enumerate}
        \item Class 1: $w_1 = \begin{bmatrix} 1 \\ 1 \end{bmatrix}$, $b_1 = 0$
        \item Class 2: $w_2 = \begin{bmatrix} 1 \\ 0 \end{bmatrix}$, $b_2 = 1$
        \item Class 3: $w_3 = \begin{bmatrix} 0 \\ 1 \end{bmatrix}$, $b_3 = -1$
    \end{enumerate}

    Draw the resulting decision boundary and clearly mark the region corresponding to each class.

    \begin{soln}

    \end{soln}
\end{prob}

\begin{prob}
    Decision trees have great expressive power: they can represent any decision boundary, by carving up the space to a fine enough granularity. Which of the other methods we've studied are similarly expressive? For each, give a brief explanation.

    \begin{subprobset}
    \begin{subprob}
        Linear classifiers

        \begin{soln}

        \end{soln}
    \end{subprob}

    \begin{subprob}
        Nearest neighbors classifiers

        \begin{soln}

        \end{soln}
    \end{subprob}

    \begin{subprob}
        Classifiers based on Gaussian generative models

        \begin{soln}

        \end{soln}
    \end{subprob}
    \end{subprobset}
\end{prob}

\begin{prob}
    In order to grow a decision tree, we need a measure of how pure or impure (in terms of labels) a node is. A popular measure of this is the Gini impurity index. If there are two labels, and a node has p fraction of one label and $1 - p$ fraction of the other, the Gini impurity index is $2p(1 - p)$. Notice that this is maximized when $p = 1/2$.

    What is the Gini impurity index for a node in which 20\% of the points have one label while 80\% have the other label?

    \begin{soln}

    \end{soln}
\end{prob}

\begin{prob}[\textit{Working with weighted data}]
    Boosting requires the weak learning algorithm to work with training data in which each point $(x^{(i)}, y^{(i)})$ has a positive weight $\lambda_i > 0$. Intuitively, a weight of two would be equivalent to having two copies of that data point. How would you incorporate weights into the following learning algorithms without explicitly making copies of data points? Give a brief explanation in each case.

    \begin{subprobset}
    \begin{subprob}
        Decision trees.

        \textit{Hint:} Previously, we measured the uncertainty at any node by looking at the fraction of points with each label (call these $p_1, \ldots, p_k$ if there are $k$ labels). Now we need to compute these fractions differently, taking weights $\lambda_i$ into account.

        \begin{soln}

        \end{soln}
    \end{subprob}

    \begin{subprob}
        Gaussian generative models.

        \textit{Hint:} For each class $j$, we need the weight of that class ($\pi_j$), the mean ($\mu_j$), and the covariance matrix ($\Sigma_j$). Now we need to compute them differently, taking weights $\lambda_i$ into account.

        \begin{soln}

        \end{soln}
    \end{subprob}

    \begin{subprob}
        Support vector machines.

        \textit{Hint:} Can you incorporate the weights $\lambda_i$ into the objective function for soft-margin SVM?

        \begin{soln}

        \end{soln}
    \end{subprob}
    \end{subprobset}
\end{prob}

\begin{prob}[\textit{Convergence behavior of boosting.}]
    Suppose we run boosting using weak classifiers from some class $H$ (such as decision stumps), and suppose that on each round, the weak learner always returns a weak classifier whose error rate on the weighted data (for that round) is at most $1/2 - \epsilon$, for some $\epsilon > 0$. Which of the following is true? Give a brief explanation in each case.

    \begin{subprobset}
    \begin{subprob}
        Boosting converges to a final classifier with zero test error.

        \begin{soln}

        \end{soln}
    \end{subprob}

    \begin{subprob}
        Boosting converges to a final classifier with zero training error.

        \begin{soln}

        \end{soln}
    \end{subprob}

    \begin{subprob}
        Boosting's final classifier belongs to class $H$.

        \begin{soln}

        \end{soln}
    \end{subprob}
    \end{subprobset}
\end{prob}

\begin{progprob}[\textit{Multiclass Perceptron}]
    Implement the multiclass Perceptron algorithm from class. Then, load in the data set \texttt{data0.txt}. This has 2-d data in four classes (coded as 0,1,2,3). Each row consists of three numbers: the two coordinates of the data points and the label.

    Run the multiclass Perceptron algorithm to learn a classifier. Turn in a plot that shows all the data points (with different colors and shapes for different labels) as well as the decision regions.

    \begin{soln}

    \end{soln}
\end{progprob}

\begin{progprob}[\textit{Multiclass SVM}]
    In this problem, we'll use support vector machines on the same data set, \texttt{data0.txt}.

    \begin{itemize}
        \item Learn a linear SVM classifier using \texttt{sklearn.svm.LinearSVC}. You will need to set \texttt{loss='hinge'} and \texttt{multi\_class='crammer\_singer'}. Try the following values of $C$: 0.01, 0.1, 1.0, 10.0.
    \end{itemize}

    \begin{subprobset}
    \begin{subprob}
        For each value of $C$, plot the decision boundary (no need to show the margins).

        \begin{soln}

        \end{soln}
    \end{subprob}

    \begin{subprob}
        What do you notice as $C$ increases? Briefly comment.

        \begin{soln}

        \end{soln}
    \end{subprob}
    \end{subprobset}
\end{progprob}

\begin{progprob}[\textit{A toy 2-d data set for decision trees and boosting}]
    Obtain the data set \texttt{mini-data.txt} from the course webpage. Each line has a two-d data point followed by a label (0 or 1).

    \begin{subprobset}
    \begin{subprob}
        Plot this data to see what it looks like. Show this plot in your writeup.

        \begin{soln}

        \end{soln}
    \end{subprob}

    \begin{subprob}
        Use \texttt{sklearn.tree.DecisionTreeClassifier} to fit a decision tree to the data. What stopping criterion did you use?

        \begin{soln}

        \end{soln}
    \end{subprob}

    \begin{subprob}
        Display the tree using \texttt{sklearn.tree.plot\_tree}.

        \begin{soln}

        \end{soln}
    \end{subprob}

    \begin{subprob}
        Fit boosted decision stumps to this data using \texttt{sklearn.ensemble.AdaBoostClassifier}. Use a relatively small number of stumps, and display each of them.

        \begin{soln}

        \end{soln}
    \end{subprob}

    \begin{subprob}
        Give a table showing how accuracy on the training data improves as each successive stump is added.

        \begin{soln}

        \end{soln}
    \end{subprob}
    \end{subprobset}
\end{progprob}
\end{probset}
\end{document}